\MajorSection{مدخل إلى الترجمة}

تضم هذه الترجمة المقدمة والقصيدة بأقسامها التسعة، ثم ملاحظتي برتون وخاتمته. نُقلت القصيدة في سطور نثرية ذات إيقاع، مع الحفاظ على ترتيب أبياتها المزدوجة وصورها ودرجة جزمها أو ترددها. ويُرقّم كل بيت مزدوج لتسهيل مراجعته مقابل النص الإنجليزي.

تنتمي مصطلحات النص ومقارناته وأرقامه إلى عالم برتون في القرن التاسع عشر. لذلك بقيت أفكاره في المادة والعقل والروح والقدر، وتصنيفاته للشعوب والأديان، ظاهرة في عباراتها. وترد الحواشي الأصلية تحت اسم «حاشية برتون»، أما التوضيحات العربية المضافة فتجتمع في آخر الكتاب تحت اسم «ملاحظات على الترجمة».

\MinorSection{الترجمة الموهومة والإطار الأدبي}

يقدّم الكتاب نفسه في صورة قصيدة للحاج عبدو اليزدي، نقلها وعلّق عليها صديقه وتلميذه «ف ب». هذا إطار أدبي وضعه برتون لعمله الإنجليزي؛ وقد احتُفظ بصوت المترجم المتخيّل في المقدمة والملاحظات والخاتمة. وكل إشارة فيها إلى «الأصل» أو إلى مقاطع «لم تُترجم» جزء من ذلك الإطار.

هذا مثال على «الترجمة الموهومة» (\textenglish{pseudotranslation}): نص يؤلّفه الكاتب بلغته، ثم يقدّمه كما لو كان ترجمة لأصل أجنبي. فبرتون ألّف القصيدة بالإنجليزية، ونسبها إلى شاعر فارسي متخيّل، وابتكر لها مترجمًا ومعلّقًا.

ومن جاذبية هذا الإطار إغراء المصدر البعيد وهيبة الحكمة المنقولة؛ فهو يكسو أفكار برتون هالة شرقية، ويستفيد من توقعات قرّاء رباعيات عمر الخيام بترجمة فيتزجيرالد. وله وظائف أخرى أيضًا: توزيع الصوت بين الشاعر والمترجم، وإجراء حوار بين وجوه الكاتب، وانتقاد حدود الفكر البريطاني المعاصر، وتوسيع مساحة الشرح والتعليق. وقد تناول غلين بورزغلوف هذه الوظائف في دراسته عن دوافع المترجم الموهوم.\footnote{انظر \textenglish{Glyn Pursglove (2011)}, ص ٤–٥؛ وترد الإحالة الكاملة في «النص المعتمد».}
